\chapter{第18套试卷——参考答案与评分标准}

\section{一、选择题答案（18分）}

\begin{enumerate}[label=\arabic*.]
\item \textbf{C} — n输入LUT本质是$2^n\times1$ SRAM，可实现任意n变量组合逻辑函数。
\item \textbf{B} — 信号在进程挂起后更新($\Delta$延时)；变量立即更新。
\item \textbf{C} — 组合进程敏感列表不完整导致锁存器。
\item \textbf{A} — Moore输出仅取决于状态；Mealy取决于状态+输入。
\item \textbf{D} — 并发语句(WHEN-ELSE, WITH-SELECT)和顺序语句(CASE, IF)均可描述组合/时序逻辑。
\item \textbf{B} — 优先级：NOT > AND > OR。\texttt{A AND NOT B OR C} = \texttt{(A AND (NOT B)) OR C}。
\end{enumerate}

\section{二、填空题答案（12分）}

\begin{enumerate}[label=\arabic*.]
\item 变量赋值\texttt{:=}立即生效；信号赋值\texttt{<=}进程挂起后生效
\item 组合进程敏感列表必须包含\textbf{所有输入信号}，否则产生\textbf{锁存器}
\item 状态机三要素：\textbf{当前状态}、\textbf{次态逻辑}、\textbf{输出逻辑}
\item 省略ELSE/WHEN OTHERS在组合进程中推断\textbf{锁存器(Latch)}
\end{enumerate}

\section{三、程序改错题解析（5分）}

\begin{enumerate}[label=错误\arabic*：]
\item PROCESS内部对SIGNAL用\texttt{:=}赋值→应为\texttt{<=}
\item IF语句缺少ELSE分支(组合进程)→推断锁存器
\item 敏感列表缺少输入信号→行为不匹配
\item 类型不匹配：STD\_LOGIC\_VECTOR与INTEGER混用→需类型转换
\item CASE缺少WHEN OTHERS→推断锁存器
\end{enumerate}

{\centering 评分：每处1分，共5分。\par}

\section{四、程序阅读分析题解析（5分）}

\textbf{（1）（1分）}环形计数器（4位，4状态one-hot循环）。\\
\textbf{（2）（1分）}初始1000→0100→0010→0001→1000，4拍一个循环。\\
\textbf{（3）（2分）}4个D触发器，one-hot编码，无需译码电路。\\
\textbf{（4）（1分）}若进入非法状态(如全0)，无法自动恢复——需加自启动逻辑。

\section{五、综合设计题参考答案（60分）}

\subsection{设计题1：8-3优先编码器（20分）}

\textbf{【VHDL】}
\begin{lstlisting}
library IEEE;
use IEEE.STD_LOGIC_1164.all;

entity priority_encoder_8to3 is
  port(I : in std_logic_vector(7 downto 0);
       Y : out std_logic_vector(2 downto 0);
       GS : out std_logic);
end priority_encoder_8to3;

architecture behav of priority_encoder_8to3 is
begin
  process(I)
  begin
    if I(7)='1' then Y<="111"; GS<='1';
    elsif I(6)='1' then Y<="110"; GS<='1';
    elsif I(5)='1' then Y<="101"; GS<='1';
    elsif I(4)='1' then Y<="100"; GS<='1';
    elsif I(3)='1' then Y<="011"; GS<='1';
    elsif I(2)='1' then Y<="010"; GS<='1';
    elsif I(1)='1' then Y<="001"; GS<='1';
    elsif I(0)='1' then Y<="000"; GS<='1';
    else Y<="000"; GS<='0'; end if;
  end process;
end behav;
\end{lstlisting}

\textbf{【评分】（20分）}ENTITY 3分 + IF-ELSIF优先级链6分 + GS信号(区分无输入)4分 + 输出编码5分 + 代码完整2分。

\subsection{设计题2：8位移位寄存器（20分）}

\textbf{【VHDL】}
\begin{lstlisting}
library IEEE;
use IEEE.STD_LOGIC_1164.all;

entity shift_reg_8bit is
  port(clk, rst : in std_logic;
       mode : in std_logic_vector(1 downto 0);
       si_l, si_r : in std_logic;
       data_in : in std_logic_vector(7 downto 0);
       q : out std_logic_vector(7 downto 0));
end shift_reg_8bit;

architecture behav of shift_reg_8bit is
  signal reg : std_logic_vector(7 downto 0);
begin
  process(clk, rst)
  begin
    if rst='1' then reg<=(others=>'0');
    elsif rising_edge(clk) then
      case mode is
        when "00" => null;                     -- 保持
        when "01" => reg <= si_r & reg(7 downto 1);  -- 右移
        when "10" => reg <= reg(6 downto 0) & si_l;  -- 左移
        when "11" => reg <= data_in;            -- 装载
        when others => null;
      end case;
    end if;
  end process;
  q <= reg;
end behav;
\end{lstlisting}

\textbf{【评分】（20分）}四模式CASE 8分 + 移位拼接正确6分 + 同步装载3分 + 异步复位2分 + 端口完整1分。

\subsection{设计题3：交通灯控制器（20分）}

\textbf{【分析】}主干道绿灯30s→黄灯5s→支路绿灯20s→黄灯5s→循环。Moore型，4状态，内部计数器计时。

\textbf{【状态定义】}
\begin{center}
\begin{tabular}{|c|l|c|c|}
\hline 状态 & 含义 & MR灯(RYG) & SR灯(RYG) \\
\hline S0 & 主干道绿灯 & 001 & 100 \\
S1 & 主干道黄灯 & 010 & 100 \\
S2 & 支路绿灯 & 100 & 001 \\
S3 & 支路黄灯 & 100 & 010 \\
\hline
\end{tabular}
\end{center}

\textbf{【VHDL核心】}
\begin{lstlisting}
library IEEE;
use IEEE.STD_LOGIC_1164.all;
use IEEE.NUMERIC_STD.all;

entity traffic_light is
  port(clk, rst : in std_logic;
       MR_light, SR_light : out std_logic_vector(2 downto 0));
end traffic_light;

architecture moore_fsm of traffic_light is
  type state_type is (MR_GREEN, MR_YELLOW, SR_GREEN, SR_YELLOW);
  signal curr, nxt : state_type;
  signal timer : integer range 0 to 29;
  constant T_GREEN : integer := 29;  -- 30秒(0-29)
  constant T_YELLOW : integer := 4;  -- 5秒(0-4)
  constant T_SGREEN : integer := 19; -- 20秒
  constant T_SYELLOW : integer := 4;
begin
  reg_proc: process(clk, rst)
  begin
    if rst='1' then curr<=MR_GREEN; timer<=0;
    elsif rising_edge(clk) then
      curr<=nxt;
      if (curr=MR_GREEN and timer=T_GREEN) or
         (curr=MR_YELLOW and timer=T_YELLOW) or
         (curr=SR_GREEN and timer=T_SGREEN) or
         (curr=SR_YELLOW and timer=T_SYELLOW) then
        timer<=0;
      else timer<=timer+1; end if;
    end if;
  end process;

  comb_proc: process(curr, timer)
  begin
    case curr is
      when MR_GREEN => if timer=T_GREEN then nxt<=MR_YELLOW; else nxt<=MR_GREEN; end if;
      when MR_YELLOW => if timer=T_YELLOW then nxt<=SR_GREEN; else nxt<=MR_YELLOW; end if;
      when SR_GREEN => if timer=T_SGREEN then nxt<=SR_YELLOW; else nxt<=SR_GREEN; end if;
      when SR_YELLOW => if timer=T_SYELLOW then nxt<=MR_GREEN; else nxt<=SR_YELLOW; end if;
    end case;
  end process;

  with curr select MR_light <= "001" when MR_GREEN, "010" when MR_YELLOW, "100" when others;
  with curr select SR_light <= "100" when MR_GREEN|MR_YELLOW, "001" when SR_GREEN, "010" when SR_YELLOW;
end moore_fsm;
\end{lstlisting}

\textbf{【评分】（20分）}4状态定义+灯控表3分 + 内部timer计时6分 + 两段式VHDL 4分 + Moore输出3分 + 时序循环正确4分。
